\section{讲座定位：理解 Mixtral 要同时记三本账}

Mixtral 常被压缩成一句话：“八个 7B 专家，每个 token 只激活两个。”这句话足够吸引注意，却不足以解释模型为什么快、为什么占显存、为什么需要复杂通信，也不足以说明所谓“专家”究竟学到了什么。本讲的正确入口是同时维护三本账：\textbf{总容量}记录所有专家能存下多少参数，\textbf{活跃计算}记录一个 token 实际经过多少参数，\textbf{系统成本}记录权重驻留、token dispatch、跨卡通信和负载不均衡。三者不能相互替代。

\subsection{先修复来源，再建立阅读路线}

本节先说明材料边界，因为旧目录存在严重来源污染：旧幻灯片文件实际是 Nathan Lambert 的 Lecture 28 对齐课程 deck，并非 Albert Jiang 的 Mixtral 讲座。CS25 官网没有给本讲发布独立 deck，因此本讲从官方 1920x1080 录像逐秒恢复 slide states，并把 34 分钟 Q\&A 作为 teacher voice 融入正文。接下来的图都来自官方录像，而不是错误 PDF。

\lecturefigure{slide-01-content.jpg}{课程路线：dense Transformer、Sparse Mixture of Experts 与 routing interpretation}{Stanford Online 官方录像 00:01:59}

\begin{importantbox}{本讲的核心问题}
Sparse Mixture of Experts（SMoE，稀疏专家混合）不是“更大的 dense model”。它把每层的单个 MLP 扩成多个 expert MLP，再由 router 为每个 token 选择少量专家。这样可以提高总容量而不同比例提高每 token FLOPs；但所有专家仍需存储，token 仍需在设备间移动，router 仍可能把负载压到少数专家。
\end{importantbox}

\begin{knowledgebox}{术语消化：三本账分别回答什么}
\begin{tabularx}{\textwidth}{L{0.18\textwidth} L{0.28\textwidth} X}
\toprule
账本 & 核心问题 & 典型变量 \\
\midrule
总容量 & 模型一共存了多少权重？ & total parameters、所有 experts、shared attention \\
活跃计算 & 一个 token 实际经过哪些权重？ & top-$k$、active parameters、per-token FLOPs \\
系统成本 & 权重和 token 怎样落在硬件上？ & HBM（High Bandwidth Memory，高带宽显存 DRAM）、expert parallelism、all-to-all、batch、load balance \\
\bottomrule
\end{tabularx}
\end{knowledgebox}

\teachervoice{00:01:04--00:01:54，Albert 把课程明确分成 architecture 与 interpretation 两部分，并反复插入 open research questions。他希望开放社区继续研究，而不是把 Mixtral 当成已经解释完的配方。}

\subsection{证据等级：哪些是测量，哪些只是工作假设}

这堂课同时包含模型配置、benchmark plot、routing histogram、社区 ablation 和 Q\&A 经验判断。阅读时必须区分：模型参数表是结构事实；2024 年 benchmark 是当时的比较证据；“MLP 存知识、attention 做推理”是 speaker 明说的 conventional wisdom；某个 expert 在数学数据上被多选只是相关性；“更大 expert count 更易专门化”则是前瞻判断。后文会在每一处标出边界。

\begin{warningbox}{名字里的 expert 不等于人类职业专家}
Expert 只是可被 router 选择的参数子网络。它可能对 token 形态、局部语法、位置、频率或多个潜在特征的线性组合敏感，并不天然对应“代码专家”“医学专家”。把 router 选择直接翻译成人类语义，是本讲要纠正的第四个 myth。
\end{warningbox}

\subsection{本章小结}

本讲的教学主线不是模型宣传，而是把 sparse capacity、active compute、resident memory、communication 与 interpretability 放进同一系统。来源修复也意味着：所有 slide coverage 和教师口头解释必须重新建立，旧讲义不能沿用。

